\section{Uniparental clocks under sole progenitorship}\label{sec:uni}

\subsection{The forced identification}
Under sole progenitorship the uniparental genealogies are forced, not
fitted: every extant mitochondrial genome descends from the founding
mother's single mtDNA molecule, and every extant Y chromosome from the
founding father's single Y. The roots of both trees \emph{are} the couple:

\begin{proposition}[Uniparental identification]\label{prop:uni}
Under an EP trajectory, the mitochondrial MRCA of all extant lineages is
the founding mother and the Y-chromosomal MRCA is the founding father; both
true dates equal $T_0$ identically, and the observed near-coincidence of
the two dates is a prediction of the model rather than a coincidence.
\end{proposition}

The second sentence deserves emphasis. In mainstream demography the
mtDNA-Eve date ($\sim$150--200 kya) and the Y-chromosomal Adam date
($\sim$120--340 kya depending on study \citep{Poznik2013,Mendez2013,
Karmin2016}) are estimates of independent coalescent functionals whose
rough alignment has no particular reason to hold; under sole
progenitorship their alignment is structural, because both equal $T_0$.
The observed overlap of the two published confidence intervals is
therefore evidence \emph{for} the single-couple structure in exactly the
sense in which the model is testable --- and it was noticed early in the
debate as a datum any origin model must accommodate
\citep{Cann1987,Vigilant1991}.

\subsection{The calibration layer}
The \emph{dates} of the two MRCAs, however, are estimates on anchored
clocks. The mtDNA rate is calibrated on the human--chimpanzee mitochondrial
divergence at an assumed split time, and the Y rate likewise. Here the
well-documented rate-calibration problem enters: the directly measured
per-generation autosomal mutation rate, $\mu \approx 1.0$--$1.3\times10^{-8}$
\citep{Scally2012,Jonsson2017}, is roughly \emph{half} the
phylogenetic rate obtained by dividing the human--chimpanzee nuclear
divergence by the paleontological split time of $\sim$6.5 Myr.
\citet{Scally2012} review the reconciliation: either ancestral generation
times were much shorter, or the ancestral mutation rate was higher, or
both --- in the standard reading, the ape-lineage mutation rate has declined
over time by a factor of roughly $2$. Any date estimated on a clock
anchored to the split therefore inherits an ancestral-rate bias: if the
true recent rate is lower than the calibration-averaged rate, recent
coalescence dates are \emph{underestimated} by the same factor.

This paper's treatment of the uniparental constraint is precisely to
apply that correction as a single clock-scaling factor $s$, with
\begin{equation}\label{eq:sdef}
\text{corrected date} \;=\; s \times \text{published estimate},
\qquad s \in [1.0,\,2.5],
\end{equation}
and to ask at which $s$ the corrected uniparental intervals contain the
joint feasibility window \eqref{eq:window}. The published estimate
landscape is taken wide: mtDNA-Eve point estimates 150--200 kya with
confidence intervals reaching 129--298 kya \citep{Vigilant1991,Cann1987};
Y-Adam estimates spanning 120--156 kya \citep{Poznik2013} through
208--338 kya \citep{Mendez2013} with intervals from 115 kya
\citep{Karmin2016} to 338 kya. Table~\ref{tab:clockscan} and
Figure~\ref{fig:clockscan} give the scan.

\begin{table}[t]
\centering\small
\caption{Uniparental clock scan: corrected mtDNA/Y intervals versus the
joint feasibility window $[400, 595]$ kya. Published mtDNA CI
\citep{Vigilant1991,Cann1987}: 129--298 kya; published Y span
\citep{Poznik2013,Mendez2013,Karmin2016}: 115--338 kya. Full scan:
\texttt{clock\_scan.csv}.}
\label{tab:clockscan}
\begin{tabular}{ccccc}
\toprule
$s$ & mtDNA CI $\times s$ (kya) & Y span $\times s$ (kya) & contains window \\
\midrule
1.00 & 129--298 & 115--338 & no \\
1.50 & 194--447 & 172--507 & no \\
1.75 & 226--522 & 201--592 & no \\
\textbf{2.00} & \textbf{258--596} & \textbf{230--676} & \textbf{yes} \\
2.25 & 290--670 & 259--760 & yes \\
2.50 & 322--745 & 288--845 & yes \\
\bottomrule
\end{tabular}
\end{table}

\begin{figure}[t]
\centering
\includegraphics[width=0.72\textwidth]{../outputs/fig7_clock_scan.pdf}
\caption{The uniparental clock scan. Corrected mtDNA (blue) and Y (red)
intervals as functions of the ancestral-rate correction factor $s$; the
orange band is the joint feasibility window
$[400, 595]$ kya. Containment holds for $s \ge 2.0$, the
reconciliation factor of \citet{Scally2012}.}
\label{fig:clockscan}
\end{figure}

The reading of Table~\ref{tab:clockscan} is the paper's principal
conditional statement, and it is stated without disguise. At the
uncorrected anchored clocks ($s = 1$), the published uniparental dates sit
a factor of roughly two below the feasibility window --- the model's
principal residual tension, stated as such. The choice between the
corrected and uncorrected readings, however, is not a choice between a
corrected model and an uncorrected one; it is a choice between two
assumptions, and only one of them is documented. Declining the correction
requires holding that a rate discrepancy demonstrated on the nuclear
genome somehow leaves the uniparental clocks untouched --- an assumption
with no published support and, on present data, no test. The asymmetry of
documented support, not the mere existence of a correction, is what
weights the layer toward the model. At $s = 2.0$ --- the
factor that \citet{Scally2012} show is required to reconcile the
pedigree-measured mutation rate with the paleontological split date
\emph{on the nuclear genome} --- the corrected mtDNA confidence interval
$[258, 596]$ kya contains the window, and the corrected Y span does so
with greater margin. The correction is not an ad hoc device introduced for
this model; it is the published reconciliation of a mainstream discrepancy,
which the standard literature applies to the split date but has not
previously applied to within-human uniparental coalescences. Applying it
uniformly is the conservative move for the EP model, and declining to apply
it is the conservative move for the standard model; the difference between
the two treatments is exactly the kind of methodological asymmetry that
formal modeling is meant to expose. Two caveats bound the layer's weight,
and both are stated rather than absorbed. First, the factor $s$ is
calibrated on \emph{nuclear} divergence --- the human--chimpanzee split
measured genome-wide --- whereas mtDNA and Y are single non-recombining
loci with calibration histories of their own; transferring the nuclear
factor to them unchanged is an assumption of the scan, adopted because it
is the only published reconciliation available and because the scan
brackets it over $s \in [1.0, 2.5]$, not because it is established for
these loci. Second, and consequently, the uniparental containment is
conditional corroboration, not a load-bearing result: every autosomal
condition of the theorem --- $\pi$, the depth range, the spectrum, LD,
TMR4A, the $\Ne$-profile --- supports the window on the direct clock with
no correction applied, and the model's feasibility stands or falls with
those alone. What is lost if the correction is refused is the uniparental
dates' agreement with the window, not the window: at $s = 1$ the published
dates prefer $T_0 \approx 200$--$300$ kya, which the diversity constraints
of Section~\ref{sec:endo} admit only at the upper edge of the primate
endowment range and with $\Ne^{\max} \ge 2\times10^4$
(Section~\ref{sec:limits}).

\section{The frequency spectrum, linkage disequilibrium, and coalescence depth}\label{sec:sfsld}

\subsection{Decomposition of the spectrum}
The EP model's site-frequency spectrum is a mixture of two components with
distinct signatures. The \emph{de-novo} component --- mutations accumulated
since $T_0$ over a population expanding from four gene copies --- is
generated by coalescent simulation (Section~\ref{sec:sims}); it carries the
strong rare-variant excess of recent growth. The \emph{backbone} component
--- the couple's endowed alleles, drifting from initial frequency
$\tfrac14$ (four-allele sites) through the expansion --- carries the
common-variant mass. At the canonical point the simulation gives
$\pi_{\mathrm{denovo}} = 0.448\times10^{-3}$ and
$\theta_{W,\mathrm{denovo}} = 1.154\times10^{-3}$ for the de-novo part,
while the backbone contributes
$\pi_{\mathrm{bb}} = H_{\mathrm{pool}}\,e^{-\Lambda_c} =
0.553\times10^{-3}$ exactly and a segregating-site density of
$1.16\times10^{-3}$ per bp (fraction $0.37$ of endowed sites still
segregating). The totals,
\begin{equation}\label{eq:sfstot}
\pi_{\mathrm{total}} = 1.001\times10^{-3}
\quad (\text{observed } 1.0\times10^{-3}), \qquad
\theta_{W,\mathrm{total}} = 1.379\times10^{-3},
\end{equation}
reproduce the observed level to within $0.2\%$; the $0.001\times10^{-3}$
gap between the simulated total and the exact master-equation value
$1.000\times10^{-3}$ of Eq.~\eqref{eq:canonval} is the de-novo
simulation's sampling noise ($0.448$ simulated against $0.447$ exact).
The excess of $\theta_W$ over $\pi$ gives
\begin{equation}\label{eq:Dval}
\text{Tajima's } D \;=\; -0.93 \quad (100\text{-kb window convention, } n = 100).
\end{equation}
Genome-wide estimates for expanding human populations are of order
$-2$ in large-sample series \citep{Tennessen2012} and of order $-1$ at
population-scale sample sizes. The EP model reproduces the sign and the
order of magnitude; it sits at the milder edge because the endowed
backbone contributes intermediate-frequency variants that a
pure-expansion model lacks, and the residual to the observed order is not
closed by this paper. The direction of the comparison, however, deserves
statement against the objection's own forecast, because the classical
diversity objection carried one: a population founded by two individuals
was expected to be depleted of rare variants, its Tajima's $D$ near zero
or positive. The EP model produces the opposite --- a strongly negative
$D$, generated by its expansion epoch on exactly the timescale the
observed spectrum demands --- because the couple carries the
common-variant backbone while mutation below $T_0$ supplies the rare
tail. The observable the objection leaned on has the sign the objection
did not expect, and the refinement that remains --- the further deepening
that population structure contributes --- moves the model toward the
data, not away from it. Two directions
are expected to account for it, and both are argued from the literature
rather than simulated here (Section~\ref{sec:limits}). Population
structure within the
post-bottleneck metapopulation --- the African subdivision that
\citet{Scerri2019} argue for on archaeological and genetic grounds ---
plausibly adds deme-private rare variants and deepens $D$ without touching
any other constraint. The deepening is argued from the literature rather
than simulated here, and until a structured simulation of the EP class is
run the spectrum condition is claimed at the level on which the objection
itself stood --- sign and order --- not detailed spectral fit. Independent SFS-based inference, separately,
recognizes a severe deep
bottleneck in human ancestry \citep{Hu2023}, confirming that the observed
spectrum is bottleneck-shaped and relocating the question from
\emph{whether} to \emph{when} --- exactly the question the EP window
\eqref{eq:window} answers with a single generation at $T_0 \approx 550$
kya. Figure~\ref{fig:sfs} shows the simulated spectrum against the
constant-size equilibrium reference.

\begin{figure}[t]
\centering
\includegraphics[width=0.72\textwidth]{../outputs/fig4_sfs.pdf}
\caption{The EP model's derived-allele-frequency spectrum (de-novo
component, coalescent simulation; backbone added separately in the text
totals) against the constant-size equilibrium reference
$\propto 1/i$. The rare-variant excess of the expansion epoch dominates
the low-frequency end.}
\label{fig:sfs}
\end{figure}

\subsection{Linkage disequilibrium}
The endowment imposes a specific initial linkage structure: at any pair of
segregating sites, the couple's four haploid genomes realize one of two
phase configurations with equal probability --- aligned phases (two
haplotypes at frequency $\tfrac12$, $r^2 = 1$) or opposite phases (four
haplotypes at $\tfrac14$, $r^2 = 0$) --- so $E[r^2_0] = \tfrac12$. The
expected decay of $E[r^2_0]$ over the post-bottleneck history is then
governed by recombination and drift. Figure~\ref{fig:ld} shows the
simulated profile (exact Wright--Fisher haplotype recursion during the
expansion, hazard-equivalent rescaled drift during the plateau,
Section~\ref{sec:sims}): $E[r^2] \approx 0.40$ at $0.1$ kb, $0.34$ at
$0.5$ kb, and $0.30$ at $1$ kb, decaying to $0.04$ by $10$ kb. The
observed genome-wide short-range values, $r^2(1\,\mathrm{kb})$ of order
$0.3$--$0.5$ across populations and minor-allele-frequency conventions
\citep{Tenesa2007}, bracket the simulated profile; the simulated value at
1 kb, $0.30$, sits at the lower edge of the observed range, and
genome-wide averages are diluted further by young de-novo pairs. The
direction of the small residual is understood: genome-wide averages
include young de-novo pairs whose LD is lower still, and population
structure raises the backbone pairs' $r^2$; the profile is
convention-sensitive at the $\pm 0.1$ level. What the simulation
establishes is that the endowment's initial two-point structure does
\emph{not} produce a detectable LD excess --- the opposite fear, that a
four-haplotype origin would imprint genome-wide haplotype blockiness, is
quantitatively false at 550 kyr of recombination time.

\begin{figure}[t]
\centering
\includegraphics[width=0.7\textwidth]{../outputs/fig5_ld.pdf}
\caption{Linkage-disequilibrium decay under the EP model at the canonical
point: $E[r^2]$ between pairs of endowed common variants against distance.
The shaded band marks the observed short-range values
$r^2(1\,\mathrm{kb}) \approx 0.3$--$0.5$ \citep{Tenesa2007}. The bimodal
structure (half the pairs start phase-aligned) produces the heavy tail
visible in the medians.}
\label{fig:ld}
\end{figure}

\subsection{Coalescence-depth distribution}
The model's expected pairwise depth, $1{,}042$ kyr at the canonical point,
is matched to the observed $\pi_{\mathrm{obs}}/2\mu$ by construction. Its
\emph{distribution}, however, is a genuine prediction: depths concentrate
just below $T_0$ for pairs coalescing in the expansion, extend to
$T_0 \approx 500$ kyr for pairs landing on the same founder copy, and
reach $T_0 + h/(2\mu) \approx 2.2$ Myr for cross-founder pairs, with
locus-to-locus variance in the endowment extending the extreme to the
$2$--$3$ Myr depths observed at the deepest autosomal loci
\citep{GarriganHammer2005,LiDurbin2011}. The mixture reproduces the
observed breadth of autosomal coalescence times ($0.5$--$2.0$ Myr typical,
deepest loci beyond $2$ Myr) without any pre-bottleneck population ---
the depth that the sampled-couple model imported from $\Nepre$ enters here
as endowed divergence. This is the formal answer to the objection that an
$N=2$ origin would leave all coalescences ``shallow'': it is the
\emph{number of lineages} at $T_0$, not their pairwise depth, that the
couple constrains, and four unrelated founder copies carry as much pairwise
divergence as a small population does.

\section{Detectability of the bottleneck}\label{sec:ne}
A natural remaining objection holds that PSMC- and MSMC-type
reconstructions \citep{LiDurbin2011,SchiffelsDurbin2014} would recover the
bottleneck. Two independent arguments say otherwise, one from the literature
and one from the present simulations. The literature argument is that a
short, sharp bottleneck leaves no trace in
reconstructions whose temporal resolution at the relevant depth is tens of
thousands of years \citep{LiDurbin2011} and whose likelihood surface is
smoothed by
regularization. The simulation argument quantifies it for the EP model:
from pairwise coalescence times under the canonical demography
(Section~\ref{sec:sims}), the hazard-inversion profile
(Figure~\ref{fig:neprof}) shows a shallow depression at the bottleneck
epoch --- the bin containing $T_0$ implies $\Ne \approx 1.2\times10^{3}$
against the plateau's $3\times10^4$ --- rather than a recovery of
$N = 2$. An order-of-magnitude depression lasting a single resolution
element, at a depth where PSMC's confidence bands are widest, is not
recoverable as a bottleneck to two; it is indistinguishable from the mild
ancient structure that real reconstructions already show. The resolution
argument is quantitative: PSMC's temporal bin width at $0.5$ Myr depth is
of order $15$--$30$ kyr \citep{LiDurbin2011}, and a bottleneck lasting a
single generation is far below it.

\begin{figure}[t]
\centering
\includegraphics[width=0.72\textwidth]{../outputs/fig6_ne_profile.pdf}
\caption{PSMC-style implied $\Ne$ from the pairwise coalescence-rate
profile under the canonical EP demography. The dashed line marks $T_0 =
550$ kya, where the true population size is $2$: the reconstruction shows
a shallow depression (implied $\Ne \approx 1.2\times10^{3}$ in the bin
containing $T_0$), not a recovery of the true size --- the bottleneck is
invisible as a bottleneck to two at PSMC resolution.}
\label{fig:neprof}
\end{figure}
